% Recreación vectorial de las doce simetrías de C2, P07.
% Los tonos distinguen sectores; la F distingue orientación y reflexión.
\begingroup
\newcommand{\hexmarca}[2]{%
\begin{tikzpicture}[scale=.42,baseline=-.6ex]
\begin{scope}[rotate=#1,xscale=#2]
\foreach \ang/\col in {0/azulAlgebra,60/verdeAlgebra,120/ocreAlgebra,180/rojoCIMAT,240/azulAlgebra!40,300/rojoCIMAT!35}{
\colorlet{sector}{\col}
\fill[sector!65] (0,0)--(\ang:1)--(\ang+60:1)--cycle;}
\draw[thick] (0:1)--(60:1)--(120:1)--(180:1)--(240:1)--(300:1)--cycle;
\draw[white,line width=2.8pt] (-.22,-.45)--(-.22,.45)--(.28,.45)
(-.22,.05)--(.18,.05);
\draw[black,line width=.8pt] (-.22,-.45)--(-.22,.45)--(.28,.45)
(-.22,.05)--(.18,.05);
\end{scope}\end{tikzpicture}}
\setlength{\tabcolsep}{8pt}
\begin{tabular}{cccccc}
\hexmarca{0}{1}&\hexmarca{60}{1}&\hexmarca{120}{1}&\hexmarca{180}{1}&\hexmarca{240}{1}&\hexmarca{300}{1}\\[7pt]
$e$&$r$&$r^2$&$r^3$&$r^4$&$r^5$\\[7pt]
\hexmarca{0}{-1}&\hexmarca{60}{-1}&\hexmarca{120}{-1}&\hexmarca{180}{-1}&\hexmarca{240}{-1}&\hexmarca{300}{-1}\\[7pt]
$s$&$rs$&$r^2s$&$r^3s$&$r^4s$&$r^5s$
\end{tabular}
\par\medskip
\begin{tikzpicture}[scale=.75,baseline=(current bounding box.center)]
\draw[thick] (0:1.2)--(60:1.2)--(120:1.2)--(180:1.2)--(240:1.2)--(300:1.2)--cycle;
\foreach \ang in {0,30,60,90,120,150}{\draw[rojoDiagramaCIMAT,dashed] (\ang:1.6)--(\ang+180:1.6);}
\draw[-{Stealth},thick] (15:1.8) arc[start angle=15,end angle=65,radius=1.8];
\node at (45:2.15) {$r$};
\end{tikzpicture}
\hspace{1cm}
\begin{tikzpicture}[scale=.6,baseline=(current bounding box.center),every node/.style={font=\small,fill=white,inner sep=2pt}]
% Ciclo de rotaciones y abanico de reflexiones presentes en la pizarra.
\foreach \k/\ang/\lab in {0/270/e,1/330/r,2/30/r^2,3/90/r^3,4/150/r^4,5/210/r^5}{
\node (R\k) at (\ang:1.8) {$\lab$};}
\foreach \a/\b in {0/1,1/2,2/3,3/4,4/5,5/0}{
\draw[-{Stealth}] (R\a)--(R\b);}
\foreach \k/\lab in {0/s,1/rs,2/r^2s,3/r^3s,4/r^4s,5/r^5s}{
\node (S\k) at ({-3.3+1.32*\k},-4) {$\lab$};
\draw (R0)--(S\k);}
\end{tikzpicture}
\endgroup
